\section{Construction of an approximate trajectory}\label{sec:approximation}

Throughout this section $L=\log(1/|t|)$, $\rho=L^{-1}$ and $\tau=s/L$.
We construct a smooth global metric by matching corrected component
metrics to exact stationary neck profiles. The endpoint matrices move
in slow time; the resulting neck profile therefore has a small time
residual, although its curvature vanishes exactly.

\subsection{Component metrics}
On a component punctured at radius $\rho/2$, put
\begin{equation}\label{eq:component-metric}
 K_i(\tau,z)=H_i(\tau)\exp(\xi_i(\tau,z)/L).
\end{equation}
It is positive Hermitian with determinant one. By
\eqref{eq:potential-bounds}, $|\xi_i|/L=O_T(\log L/L)$ there.
Transfer its matrix entries to the corresponding part of $X_t$
using Lemma~\ref{lem:plumbing}.

\begin{lemma}\label{lem:component-residual}
On the transferred component regions,
\begin{equation}\label{eq:component-residual}
 \int |R(K_i)|\,dA_t
 \le C_T\frac{1+\log L}{L^2}+C_Te^{-L}L^n.
\end{equation}
On every fixed compact subset away from the punctures, the residual
and its fixed finite spatial derivatives are $O_T(L^{-2})$.
\end{lemma}
\begin{proof}
First calculate on the central component. Write $h=\xi_i/L$ and
suppress the component index. For bounded small $h$,
\[
 e^{-h}\partial e^h
 =\int_0^1 e^{-ah}(\partial h)e^{ah}\,da
 =\Phi(\operatorname{ad}_h)\partial h,
\]
where $\Phi$ is smooth, $\Phi(0)=I$, and its derivatives are bounded.
After applying $\bar\partial$ and contracting, the difference from
the linear term is bounded by
\[
 |2i\Lambda F_{He^h}-\Delta h|
 \le C_T(|h||\Delta h|+|\nabla h|^2).
\]
All commutators remain in this expression. The reference matrix $H$
is spatially constant. For the time derivative,
\[
 (He^h)^{-1}\partial_s(He^h)
 =L^{-1}e^{-h}Ve^h+e^{-h}\partial_s e^h,
\]
whose difference from $V/L$ is bounded by
$C_TL^{-2}(|\xi||V|+|\partial_\tau\xi|)$. Thus
\begin{equation}\label{eq:thick-expansion}
 R(K_i)=L^{-1}(V_i+\Delta_i\xi_i)
 +O_T\!\left(L^{-2}\bigl(|\xi_i||V_i|+|\partial_\tau\xi_i|
             +|\nabla\xi_i|^2+|\xi_i||\Delta_i\xi_i|\bigr)\right).
\end{equation}
The first term vanishes by \eqref{eq:compatibility}. Integration of
the other terms uses \eqref{eq:potential-bounds} and proves the first
term of \eqref{eq:component-residual}. The coordinate and coefficient
errors in Lemma~\ref{lem:plumbing}, multiplied by the polynomial
derivative bounds of the potentials, contribute $O_T(e^{-L}L^n)$.
On a fixed compact subset every potential derivative is bounded,
which proves the local assertion.
\end{proof}

\subsection{Geodesic necks}
At a node $p$ joining $C_i,C_j$, write
\begin{equation}\label{eq:regular-parts}
 \xi_i=-B_{ij}\log|u|+c_{i,p}+q_{i,p}(u),\qquad q_{i,p}(0)=0,
\end{equation}
where the regular part is smooth, with bounded slow-time derivatives.
Use the analogous expansion on the $v$ branch. Set
\[
 P_{i,p}=H_i\exp(c_{i,p}/L),\qquad
 P_{j,p}=H_j\exp(c_{j,p}/L),
\]
and on $uv=t$ define
\begin{equation}\label{eq:neck-profile}
 N_p(\tau,u)=\Geo(P_{i,p},P_{j,p};\theta),
 \qquad \theta=-\frac{\log|u|}{L}.
\end{equation}
Since $1-\theta=-\log|v|/L$, the same expression from the opposite
end is obtained by reversing the geodesic. In the neck region used
below, $0\le\theta\le1$.

\begin{lemma}\label{lem:neck}
For every fixed $\tau$, the metric $N_p$ has zero Chern curvature.
On $|u|,|v|\le2\rho$ its time residual satisfies
\begin{equation}\label{eq:neck-residual}
 |R(N_p)|\le C_T/L,\qquad
 \int |R(N_p)|\,dA_t\le C_T\rho^2/L.
\end{equation}
\end{lemma}
\begin{proof}
Put $S=\log(P_{i,p}^{-1/2}P_{j,p}P_{i,p}^{-1/2})$. The spatial
connection is exactly
\[
 N_p^{-1}\partial N_p
 =P_{i,p}^{-1/2}SP_{i,p}^{1/2}\,\partial\theta.
\]
The coefficient matrix is spatially constant and $\theta$ is harmonic,
so its $\bar\partial$ derivative is zero. The matrices $P_{i,p}$ and
their slow-time derivatives stay in bounded sets, and geodesic
interpolation is smooth on their compact range. Thus
$|N_p^{-1}\partial_sN_p|\le C_T/L$. The area estimate is
Lemma~\ref{lem:plumbing}(3).
\end{proof}

\subsection{Differentiated matching}
Near the $i$ end use the logarithmic matrix coordinate
$\ell_i(K)=\log(H_i^{-1}K)$. It takes values in the space of
$H_i$-self-adjoint trace-free endomorphisms. All its derivatives
needed here are bounded on the compact matrix range.

\begin{lemma}\label{lem:matching}
On the transition annulus $\rho\le|u|\le2\rho$, for $k=0,1,2$,
\begin{equation}\label{eq:matching}
 |D^k(\ell_i(K_i)-\ell_i(N_p))|
 \le C_T\left(\frac{\rho^{1-k}}{L}
       +\frac{1+|\log\rho|}{L^2\rho^k}\right).
\end{equation}
The same estimate holds after any fixed number of slow-time
derivatives. The corresponding assertion holds at the other end.
\end{lemma}
\begin{proof}
Write $\varepsilon=L^{-1}$ and, at a fixed slow time, consider
\[
 F(\varepsilon,\theta)=\ell_i\!\left(
 \Geo(H_i e^{\varepsilon c_{i,p}},
       H_j e^{\varepsilon c_{j,p}};\theta)\right).
\]
Its derivatives are uniformly bounded for $\theta\in[0,1]$ and
small $\varepsilon$. Exactly,
$F(0,\theta)=\theta B_{ij}$ and $F(\varepsilon,0)=\varepsilon c_{i,p}$.
The difference from their sum vanishes on both coordinate axes.
The double integral of its mixed derivative therefore gives
\[
 F(\varepsilon,\theta)
 =\theta B_{ij}+\varepsilon c_{i,p}
     +\varepsilon\theta G(\varepsilon,\theta,\tau),
\]
with $G$ and its required derivatives uniformly bounded. On the
annulus, $\theta=-\varepsilon\log|u|$ and
$D^k\theta=O(\varepsilon\rho^{-k})$ for $k=1,2$. The component
logarithm is
\[
 \ell_i(K_i)=\theta B_{ij}+\varepsilon c_{i,p}
                         +\varepsilon q_{i,p}(u).
\]
Since $q_{i,p}$ is smooth and zero at the puncture,
$|D^kq_{i,p}|\le C_T\rho^{1-k}$ for $k\le2$, using a harmless
weaker bound when $k=2$. Subtraction and the chain rule give
\eqref{eq:matching}. The same calculation after slow-time
differentiation uses the bounded derivatives of the finite flow.
No pair of the matrices has been assumed to commute.
\end{proof}

Choose a cutoff $\chi(|u|/\rho)$ equal to zero near the inner end
and one near the outer end of the transition annulus. Set there
\begin{equation}\label{eq:interpolation}
 \widehat H_t=H_i\exp\!\left(
       \ell_i(N_p)+\chi(\ell_i(K_i)-\ell_i(N_p))\right).
\end{equation}
Use $K_i$ on the thick pieces and $N_p$ on the middle necks.
The formula is positive Hermitian with determinant one and agrees
with the adjacent metrics on neighbourhoods of the boundaries.
It thus defines a smooth global metric. Positivity and uniform
two-sided bounds follow from compactness of the component tuple,
smallness of $\xi_i/L$ on the thick pieces, and compactness of the
neck geodesics.

\begin{proposition}\label{prop:residual}
For every fixed $T<\infty$, the constructed metrics satisfy
\begin{equation}\label{eq:residual-spatial}
 \sup_{0\le s\le TL}\int_{X_t}|R(\widehat H_t(s))|\,dA_t
 \le C_T\frac{(1+\log L)^2}{L^2},
\end{equation}
and consequently
\begin{equation}\label{eq:residual-integrated}
 \int_0^{TL}\!\int_{X_t}|R(\widehat H_t(s))|\,dA_t\,ds
 \le C_T\frac{(1+\log L)^2}{L}.
\end{equation}
On a fixed compact component subset the residual is $O_T(L^{-2})$,
and \eqref{eq:approx-thick-intro} holds.
\end{proposition}
\begin{proof}
It remains to estimate the transition annuli. Put
$n=\ell_i(N_p)$ and $\delta=\chi(\ell_i(K_i)-n)$.
The curvature operator in logarithmic coordinates is a second-order
expression with smooth coefficients on the compact matrix range.
Compare its value at $n+\delta$ with its value at $n$, which is
exactly zero. The difference is bounded by a constant times
\begin{equation}\label{eq:operator-difference}
 |D^2\delta|+|Dn||D\delta|
 +|\delta|(|D^2n|+|Dn|^2)+|D\delta|^2.
\end{equation}
Here $|Dn|\le C_T/(L\rho)$,
$|D^2n|\le C_T/(L\rho^2)$, and each cutoff derivative of order
$k$ is $O(\rho^{-k})$. Applying \eqref{eq:matching} in
\eqref{eq:operator-difference} and integrating over area
$O(\rho^2)$ gives
\[
 \int_{\rm transition}|2i\Lambda F_{\widehat H_t}|\,dA_t
 \le C_T\left(\frac\rho L+\frac{(1+\log L)^2}{L^2}\right).
\]
The logarithmic time derivative is bounded by $C_T/L$ throughout
the annulus, so its integral is at most $C_T\rho^2/L$.
In particular it fits the same bound. Using the exactly flat
profile in \eqref{eq:operator-difference} is essential: a bound on
the full interpolant Hessian would produce an order $L^{-1}$
integral and would not suffice after time integration.

Sum these estimates over finitely many nodes and components, and
use Lemmas~\ref{lem:component-residual} and~\ref{lem:neck}.
For $\rho=L^{-1}$ the exponentially small plumbing errors are
absorbed in the right side of \eqref{eq:residual-spatial}.
Integration over the physical interval of length $TL$ proves
\eqref{eq:residual-integrated}. A fixed compact component subset
eventually lies entirely in a thick region. There the last assertion
follows from \eqref{eq:component-metric} and
Lemma~\ref{lem:component-residual}.
\end{proof}
